\section{Conclusion}

We opened by noting that 78\% of benchmark selection effort is automatable, yet researchers still spend 2--4 weeks on manual review. Our work demonstrates this automation is achievable through systematic extraction of benchmark design features and historical validation of coverage prediction.

\subsection{Summary}

We addressed the benchmark selection bottleneck by discovering that \textbf{modality, not task type, drives coverage constraints}. Design features (task formulation, evaluation metrics, data modality, dataset characteristics) can be extracted objectively (Cohen's kappa $\geq 0.917$) and clustered into coverage families that predict future citation co-occurrence with 78\% accuracy---validated via historical train/test split preventing circular reasoning.

Our main contributions are:

\begin{enumerate}
\item \textbf{Temporal persistence demonstration:} Pre-2023 features predict 2023--2024 patterns with 78.07\% citation overlap, outperforming random (19.61\%) by 585\%.

\item \textbf{Standardized extraction achieving substantial agreement:} Kappa 0.917--1.0 enables scaling from pilot (20 benchmarks) to large-scale analysis (100+).

\item \textbf{Modality-driven clustering discovery:} 0.748 average intra-family similarity reveals modality creates stronger coverage patterns than task formulation, challenging common assumptions.
\end{enumerate}

\subsection{Future Directions}

This work opens several promising avenues:

\textbf{From untested alternatives:} Real-world citation validation on 1000+ ArXiv papers to measure precision degradation from synthetic (100\%) to realistic (expected 75--85\%). Task-based clustering ablation to isolate modality vs task contributions.

\textbf{From scope extensions:} Scale to 100+ benchmarks covering rare modalities (video, 3D, tabular) to discover additional families. Fine-grained subclusters ($k=8$--12) to capture task-based patterns within modality groups.

\textbf{From unverified assumptions:} Cross-temporal robustness testing (pre-2020 $\rightarrow$ 2024) to assess 4-year persistence. Citation threshold sensitivity analysis (10/25/50/100 citations) to determine minimum coverage requirements.

Beyond scaling, this work enables \textbf{automated coverage gap discovery}---identifying hypothesis categories with zero existing benchmarks before implementation begins---and integration into research planning tools to estimate validation feasibility during hypothesis formulation rather than after weeks of failed benchmark searches.

Modality, not task type, drives benchmark coverage---a simple insight that enables systematic prediction where only manual review existed before, reducing weeks to minutes.
